% Generated by report/make_restated.py from theory/A1_lower_bound.tex: do not edit (edit the note and rerun).
\begin{restated}{Theorem}{thm:HLOY}{A1}{0}{Hannah--Liu--O'Connor--Yin bounds; \cite{hannah2018breaking}}
Let the components be $L$-smooth and $F$ be $\mu$-strongly convex, $\kappa=L/\mu$.
\begin{enumerate}[nosep,label=(\alph*)]
\item (Their Corollary~3, a rewording of Theorem~\ref{thm:LZ}.) For every $\varepsilon$ and every randomized algorithm following the span assumption there are a dimension $d$ and $L$-smooth, $\mu$-strongly convex $f_i$ on $\R^d$ such that reaching $\E f(x^k)-f(x^\ast)<\varepsilon$ takes the algorithm at least $\Omega((n+\sqrt{\kappa n})\ln(1/\varepsilon))$ steps.
\item (Upper bound; their Corollary~2.) With epoch length $m=n+121\kappa_Q$ and step size $\eta=\kappa_Q^{1/2}m^{-1/2}/(2L_Q)$, Prox-SVRG needs
\[
K(\varepsilon)=O\Big(\Big(\frac{n}{1+(\ln(n/\kappa_Q))_+}+\kappa_Q\Big)\ln\frac1\varepsilon+n+\kappa_Q\Big)
\]
gradient evaluations to reach $\E f(x^{K})-f(x^\ast)<\varepsilon$ (here $(a)_+=\max\{a,0\}$ and $\kappa_Q=L_Q/\mu$).
\item (Lower bound; their Theorem~2, stated in full as their Theorem~4.) For any oblivious $p$-CLI algorithm (their Appendix~B lists ``SDCA, SAG, SAGA, SVRG, SARAH, etc.\ without proximal terms'', and they state before their Theorem~2 that the class ``contains all VR algorithms mentioned in this paper'', Katyusha included; L-Katyusha with fixed parameters, which is posterior to \cite{hannah2018breaking}, is an oblivious CLI as well, since its coin flips and parameters do not depend on the oracle answers: this last inclusion is ours) and all $\mu,L$ there exist $L$-smooth, $\mu$-strongly convex $f_i$ such that at least
\[
K(\varepsilon)=\tilde\Omega\Big(\Big(\frac{n}{1+(\ln(n/\kappa))_+}+\sqrt{n\kappa}\Big)\ln\frac1\varepsilon+n\Big)
\]
iterations are needed; equivalently (their (1.3)) $K(\varepsilon)=\Omega(n+\sqrt{n\kappa}\ln(1/\varepsilon))$ for $n=O(\kappa)$ and $K(\varepsilon)=\Omega\big(n+\frac{n}{1+(\ln(n/\kappa))_+}\ln(1/\varepsilon)\big)$ for $\kappa=O(n)$.
\item (Example; their Corollary~1.) For $\varepsilon=1/n^\alpha$ and $\kappa=n^\beta$ with $\alpha,\beta\in(0,1)$, modified SVRG needs $O(n)$ iterations whereas algorithms following the span assumption need $\Omega(n\ln n)$ (the source writes ``require $O(n\ln(n))$ iterations'', quoting Theorem~\ref{thm:LZ}; the content is a lower bound).
\end{enumerate}
\end{restated}
